\documentclass[11pt]{article}
\usepackage[margin=1in]{geometry}
\usepackage{amsmath,amssymb,amsthm,bm,enumitem,natbib,hyperref,cleveref}
\newtheorem{theorem}{Theorem}
\newtheorem{lemma}{Lemma}
\newtheorem{proposition}{Proposition}
\newtheorem{assumption}{Assumption}
\crefname{assumption}{Assumption}{Assumptions}
\Crefname{assumption}{Assumption}{Assumptions}
\newcommand{\E}{\mathbb E}
\newcommand{\initscale}{\varepsilon}
\newcommand{\rulecredit}{\Phi_{\mathrm{rule}}}
\newcommand{\attconst}{C_{\mathrm{att}}}
\newcommand{\pd}[2]{\frac{\partial #1}{\partial #2}}
\title{Theorem 2: Conditional Gaussian Rule-Credit Addendum}
\date{}
\setlength{\emergencystretch}{1em}
\begin{document}
\maketitle
This standalone preview displays only the new part (iii) and its proof.
The manuscript retains its original parts (i) and (ii).
\begin{assumption}[Architecture]\label{ass:arch}
The manuscript's two-block causal Transformer, with full-prefix attention and a linear readout.
\end{assumption}
\begin{assumption}[Geometric initialization]\label{ass:init}
Normalized matrix norms, value-embedding conditioning and nonzero MLP bias magnitudes obey deterministic envelopes. The initialization scale is below the resulting deterministic threshold.
\end{assumption}
\begin{assumption}[Task]\label{ass:task}
Modular counting with a fixed query, independent letters, and a uniform start.
\end{assumption}
\begin{lemma}[Leading-coefficient input from the manuscript]\label{lem:leading}
\begin{enumerate}[label=(\roman*)]
\item The selected readout-gradient coordinate is $F_n(\varepsilon)=\varepsilon^6P_n+r_n$,
where $P_n$ is homogeneous of degree six and
$|r_n|\le(C_{\mathrm{rem}}/n)(\varepsilon^7/\varepsilon_0+\varepsilon^{10})$
on the geometric event.
\item The distinguished monomial has coefficient at least
$p(1-p)\min\{1-2p,1/4\}/(3Kn\sqrt{d_{\mathrm h}})$.
Its entries are $\tilde E_{s^\ast,0}$, $\tilde E_{w^\star,1}$,
$\tilde E_{w^\star,2}$ and one entry each of the selected block-2 head's
query, key and value matrices: $(\tilde W_Q)_{0,0}$, $(\tilde W_K)_{0,1}$,
$(\tilde W_V)_{b,2}$, with readout row $a=s^\ast+1$.
\end{enumerate}
\end{lemma}
\begin{theorem}[New part (iii)]
\begin{enumerate}[label=(\roman*),start=3]
\input{theorem2_main_insert}
\end{enumerate}
\end{theorem}
% ============================================================================
% Appendix insert: place immediately after the proof of Lemma C.8, before the
% existing ``Proof of Theorem 2(i)'' (or after it if preferred).  Needs amsmath,
% amssymb, amsthm; a lemma environment and proof environment must exist.
% Also needs bibliography items for Carbery--Wright and Gaussian
% hypercontractivity (see theorem2_references.bib).
% ============================================================================

\subsubsection*{Quantitative rule credit under Gaussian initialization}
\label{app:t2-gaussian}

We refine \cref{lem:leading} from generic nonvanishing to a quantitative
high-probability bound.  Write
\begin{equation}
 F_n(\varepsilon)=
 \left.\frac{\partial\Phi_{\mathrm{rule}}}{\partial V_{ab}}
 \right|_{\theta=\varepsilon\widetilde\theta}
 =\varepsilon^6P_n(\widetilde\theta)+r_n(\varepsilon),
 \qquad
 |r_n(\varepsilon)|\le\frac{C_{\mathrm{rem}}}{n}
 \left(\frac{\varepsilon^7}{\varepsilon_0}+\varepsilon^{10}\right).
 \label{eq:quant-rule-expansion}
\end{equation}
We may weaken this to $C_{\mathrm{rem}}(\varepsilon^7/(n\varepsilon_0)+\varepsilon^{10})$
when using \cref{eq:t2-gaussian-epsilon}. The remainder estimate holds whenever the deterministic hypotheses of
\cref{lem:leading} are satisfied and
$\varepsilon\le\varepsilon_0$,
$C_{\mathrm{att}}\varepsilon^4\le1/2$.

\begin{proposition}[Conditional Gaussian rule credit]
\label{lem:quant-rule-gradient}
Under \cref{ass:arch,ass:task} with full-prefix attention, uniform start and
$0<p<1/2$, suppose the normalized attention/embedding coordinates of $P_n$
are independent centered Gaussians. Its six distinguished coordinates have
standard deviations $\sigma_1,\ldots,\sigma_6>0$. Fix $0<\delta<1$.
Assume an event $\mathcal E_{\mathrm{geom}}$, of probability at least
$1-\delta/2$, on which \cref{ass:init} and the expansion of
\cref{lem:leading} hold with deterministic envelopes
$\varepsilon_0,C_{\mathrm{att}},C_{\mathrm{rem}}>0$. Define
\begin{equation}\label{eq:t2-gaussian-threshold}
 A_{n,\delta}=\frac{p(1-p)\min\{1-2p,1/4\}}{3Kn\sqrt{d_{\mathrm h}}}
 \left(\prod_{i=1}^6\sigma_i\right)
 \left(\frac{\delta}{2C_6}\right)^6,
\end{equation}
where $C_6\ge1$ is an absolute Gaussian anti-concentration constant.
For any deterministic initialization scale satisfying
\begin{equation}\label{eq:t2-gaussian-epsilon}
 \varepsilon\le\min\left\{\varepsilon_0,(2C_{\mathrm{att}})^{-1/4},
 \frac{n\varepsilon_0 A_{n,\delta}}{4C_{\mathrm{rem}}},
 \left(\frac{A_{n,\delta}}{4C_{\mathrm{rem}}}\right)^{1/4}\right\},
\end{equation}
\begin{equation}\label{eq:t2-gaussian-credit}
 \Pr\!\left\{\|\nabla_{\bm\theta}\rulecredit
 (\varepsilon\tilde{\bm\theta})\|
 \ge\tfrac12\varepsilon^6 A_{n,\delta}\right\}\ge1-\delta.
\end{equation}
Equivalently, the main-text constant is
$c_{\mathrm G}(p,K)=p(1-p)\min\{1-2p,1/4\}/[6K(2C_6)^6]$.

\end{proposition}
\begin{proof}
\emph{Gaussian projection.} Let $X_1,\ldots,X_6$ be the distinct coordinates
identified in \cref{lem:leading}(ii), and set $Y=\prod_i(X_i/\sigma_i)$.
Then $\mathbb E Y^2=1$. A monomial has nonzero expectation against $Y$
only if it contains each selected variable to a positive odd power.
Total degree six forces exactly their square-free product; hence
$\mathbb E[P_nY]=c_n\prod_i\sigma_i$. Cauchy--Schwarz and the existing coefficient bound give
\begin{equation}\label{eq:quant-rule-L2}
 \|P_n\|_{L^2}\ge |c_n|\prod_i\sigma_i
 \ge\frac{p(1-p)\min\{1-2p,1/4\}}{3Kn\sqrt{d_h}}\prod_i\sigma_i.
\end{equation}
\emph{Anti-concentration.} Rescale the independent coordinates to standard
Gaussians. Carbery--Wright gives
$\Pr(|P_n|\le t)\le6C_{\mathrm{CW}}(t/\|P_n\|_{L^1})^{1/6}$
\citep{carbery2001distributional}. Gaussian hypercontractivity gives
$\|P_n\|_{L^4}\le27\|P_n\|_{L^2}$ \citep{janson1997gaussian}; interpolation
$\|P_n\|_{L^2}\le\|P_n\|_{L^1}^{1/3}\|P_n\|_{L^4}^{2/3}$ implies
$\|P_n\|_{L^1}\ge3^{-6}\|P_n\|_{L^2}$.
Thus, with $C_6=\max\{1,18C_{\mathrm{CW}}\}$ and $0<\rho<1$,
\begin{equation}\label{eq:quant-rule-AC}
 \Pr\!\left\{|P_n|\ge(\rho/C_6)^6\|P_n\|_{L^2}\right\}\ge1-\rho.
\end{equation}
\emph{Gradient transfer.} Take $\rho=\delta/2$. The signal event
$|P_n|\ge A_{n,\delta}$ and the geometry event each fail with probability
at most $\delta/2$, so their intersection has probability at least
$1-\delta$, without any independence assumption. On that intersection,
\begin{align*}
 |r_n(\varepsilon)|/\varepsilon^6
 &\le C_{\mathrm{rem}}(\varepsilon/(n\varepsilon_0)+\varepsilon^4)
 \le A_{n,\delta}/2,\\
 \|\nabla_\theta\Phi_{\mathrm{rule}}\|
 &\ge |F_n(\varepsilon)|\ge\varepsilon^6|P_n|-|r_n(\varepsilon)|
 \ge\tfrac12\varepsilon^6 A_{n,\delta}.
\end{align*}
The scale conditions make each normalized remainder term at most
$A_{n,\delta}/4$.
\end{proof}



For $K=2,p=1/4$, the coefficient lower bound is $1/(128n\sqrt{d_{\mathrm h}})$. Inverse-polynomial scale/credit
requires corresponding control of $\sigma_i^{-1}$, $\varepsilon_0^{-1}$
and the envelope-dependent expansion constants; the geometry event alone
does not establish these uniform bounds.

\paragraph{An explicit Gaussian-weight geometry event.}
\label{app:t2-gaussian-geometry}
The conditional event can be verified for the following specified law.
Fix all widths with $d_{\mathrm{res}}\ge\max\{K,3\}$. Draw each normalized
matrix/embedding block $M_b\in\mathbb R^{r_b\times c_b}$ independently
with independent entries $N(0,\tau_b^2)$, where every $\tau_b>0$ is
specified before sampling. A zero readout may instead be fixed. Let $B$
be the number of Gaussian blocks, and let value-token embedding rows
have common entry standard deviation $\tau_E$. Set
 $a_0=\max_b 2\tau_b\sqrt{2[(r_b+c_b)\log9+\log(8B/\delta)]},\qquad c_{\mathrm{emb}}=\frac{\delta\tau_E}{4K^{3/2}}.$ 
Use normalized MLP biases with $1\le|\tilde b_m|\le2$ and fix
$N_b=2\sqrt{w_{\max}}$, where $w_{\max}$ is the largest MLP width.
All relative biases are fixed initially at zero. Then, with probability
at least $1-\delta/2$, every block has operator norm at most $a_0$ and
$\sigma_{\min}(E_{\mathrm{val}})\ge c_{\mathrm{emb}}$. Thus the geometric
conditions of \cref{ass:init} have fixed, nonrandom envelopes; choose
$\varepsilon_0$ below its threshold and those of \cref{lem:leading}.
The resulting $C_{\mathrm{att}},C_{\mathrm{rem}}$ are the deterministic
constants for these envelopes, not values fitted after observing a seed.

For the operator-norm assertion, take $1/4$-nets on both unit spheres,
each of cardinality at most $9^{\mathrm{dimension}}$. The norm is at most
twice the maximal bilinear form on these nets. Gaussian tails and a union
bound give failure probability at most $\delta/(4B)$ for each displayed
block bound. For embedding conditioning, let $D_i$ be the distance of
row $i$ to the span of the other $K-1$ rows. Conditional on those rows,
at least one orthogonal projection is $N(0,\tau_E^2)$, so
$\Pr(D_i<t)\le\sqrt{2/\pi}\,t/\tau_E$. For any unit combination of rows,
choose a coefficient of magnitude at least $K^{-1/2}$ and project
orthogonally to the span of the other rows; hence
$\sigma_{\min}(E_{\mathrm{val}})\ge\min_i D_i/\sqrt K$.
Taking $t=\delta\tau_E/(4K)$ bounds conditioning failure by $\delta/4$.
The two failure bounds sum to at most $\delta/2$.

Fan-in Gaussian weights are covered by taking $\tau_b=c_b^{-1/2}$ and
embedding variance $\tau_E^2=d_{\mathrm{res}}^{-1}$. This explicit witness
still uses non-Gaussian, bounded-away-from-zero biases. It does not cover
ordinary zero-bias or fully Gaussian-bias initialization. Nor does it
assert that all envelope-dependent expansion constants are width-uniform.
For the formal-regime bias law $\tilde b_m=\operatorname{sign}(Z_m)(1+|Z_m|)$
in \nolinkurl{handcoded/regime_both.py}, with independent standard-normal
$Z_m$ and $N_{\mathrm{bias}}$ such coordinates, instead use
 $N_b=\sqrt{w_{\max}}\left(1+\sqrt{2\log(16N_{\mathrm{bias}}/\delta)}\right),$ 
and replace $8B/\delta$ above by $16B/\delta$. Matrix norms, bias norms,
and embedding conditioning then fail with total probability at most
$\delta/8+\delta/8+\delta/4=\delta/2$.

\bibliographystyle{plainnat}
\bibliography{theorem2_references}
\end{document}
